\documentclass{book}
\usepackage[UTF8]{ctex}
\usepackage{geometry,amsmath,amssymb,hyperref,xcolor,latexsym}
\geometry{portrait}

%%%%%%%%%% Start TeXmacs macros
\newcommand{\assign}{:=}
\newcommand{\mathd}{\mathrm{d}}
\newcommand{\mathi}{\mathrm{i}}
\newcommand{\mathpi}{\pi}
\newtheorem{axiom}{公理}[chapter]
\newtheorem{corollary}{推论}[chapter]
\newtheorem{lemma}{引理}[chapter]
\newtheorem{definition}{定义}[chapter]
\newtheorem{remark}{注}[chapter]
\newcommand{\todo}[1]{{\color{red}#1}}
%%%%%%%%%% End TeXmacs macros

\begin{document}

\title{一般流形上的路径积分}
\author{}
\date{\today}

\maketitle

\chapter{引言}

在文献 \emph{从零构建量子力学}（quantum.tex）中，我们从三条基本公理出发推导了量子力学的时间演化方程。推导的关键步骤是路径积分形式，其中我们用到了 $\mathbb{R}^d$ 上的傅里叶变换来得到 $k\dot{x}$ 项：
\[
\delta(x-y) = \int_{\mathbb{R}^d} \frac{\mathd k}{(2\mathpi)^d} e^{\mathi k (x-y)},
\]
以及核 $r(x,y,t)$ 的类似表示。这个操作依赖于 $\mathbb{R}^d$ 的\textbf{平移结构}——表达式 $x-y$ 只有在向量空间中才有意义。

但是，量子力学并不局限于平直空间。在广义相对论、凝聚态物理中的弯曲晶格、以及量子引力等场景中，我们需要将量子力学推广到一般的微分流形 $M$ 上。这就引出了一个基本问题：

\begin{em}
  路径积分能否从 $\mathbb{R}^d$ 推广到一般的光滑流形 $M$？
\end{em}

直接照搬 $\mathbb{R}^d$ 的方法会遇到根本困难：$M$ 不是向量空间，$x-y$ 没有定义，因此 $\delta(x-y)$ 的傅里叶表示不再全局有效。然而，$M$ 在每一点 $y$ 处都有一个切空间 $T_yM$（它是一个 $d$ 维向量空间），这为我们提供了局部替代平移结构的工具。

本篇的核心策略是：

\begin{enumerate}
  \item 用\textbf{指数映射}（exponential map）将 $M$ 上的点拉回到切空间上操作；
  \item 在切空间上使用傅里叶变换；
  \item 证明所得的作用量 $S = \int [p\dot{x} - H(p,x)]\mathd t$ 中的 $p\dot{x}$ 项是余切丛上的\textbf{典范 1-形式}，具有坐标无关性；
  \item 流形上的对称性（一般协变性，即 $\operatorname{Diff}(M)$ 不变性）强制截断阶 $N_{\mathrm{cut}} = 2$，得到 Laplace--Beltrami 算子作为动能项。
\end{enumerate}

我们将在以下各章逐步展开这一方案。

\chapter{Riemann 几何基础}

在进入物理之前，我们需要复习或建立必要的微分几何工具。这些概念将贯穿全文。

\section{微分流形与切空间}

\begin{definition}
  一个 $d$ 维光滑流形 $M$ 是一个 Hausdorff 拓扑空间，其每一点 $y \in M$ 都有一个邻域 $U_y$ 同胚于 $\mathbb{R}^d$ 中的开集。这个同胚映射 $\psi_y: U_y \to \mathbb{R}^d$ 称为\textbf{坐标图}（chart），其逆 $\psi_y^{-1}: \mathbb{R}^d \to U_y$ 称为\textbf{参数化}（parametrization）。
\end{definition}

流形的关键特征是：它在局部看起来像 $\mathbb{R}^d$，但整体可能有复杂的拓扑结构（如球面、环面等）。

\begin{definition}
  在点 $y \in M$ 处的\textbf{切向量} $\xi$ 是经过 $y$ 的光滑曲线 $\gamma: (-1,1) \to M$（满足 $\gamma(0)=y$）的等价类，其中两条曲线等价当且仅当它们在某个坐标图中的导数相同。所有切向量的集合构成 $d$ 维向量空间 $T_yM$，称为 $y$ 处的\textbf{切空间}。
\end{definition}

在坐标图 $(U_y, \psi_y)$ 中，取 $\psi_y = (x^1, \ldots, x^d)$，则切空间 $T_yM$ 有自然基 $\{ \partial/\partial x^\alpha |_y \}$，其中 $\alpha=1,\ldots,d$。任意切向量 $\xi \in T_yM$ 可写为 $\xi = \xi^\alpha \partial/\partial x^\alpha$。

\begin{definition}
  $T_yM$ 的对偶空间称为\textbf{余切空间} $T_y^*M$，其中的元素 $p$ 称为\textbf{余切向量}。$p$ 是 $T_yM$ 上的线性泛函：$p: T_yM \to \mathbb{R}$。在坐标基下，$p = p_\alpha \mathd x^\alpha$，且配对（pairing）为
  \[
  p(\xi) = p_\alpha \xi^\alpha,
  \]
  其中 $p_\alpha$ 是协变分量，$\xi^\alpha$ 是逆变分量。重复指标遵循 Einstein 求和约定。
\end{definition}

整个流形上所有切空间的并集 $TM = \bigcup_{y\in M} T_yM$ 称为\textbf{切丛}，类似地 $T^*M = \bigcup_{y\in M} T_y^*M$ 称为\textbf{余切丛}。

\section{Riemann 度规与体积元}

为了定义长度、角度和体积，流形上需要额外的结构。

\begin{definition}
  一个\textbf{Riemann 度规} $g$ 是 $M$ 上光滑的 $(0,2)$ 型对称正定张量场。在局部坐标 $(x^1,\ldots,x^d)$ 中，
  \[
  g = g_{\alpha\beta}(x) \mathd x^\alpha \otimes \mathd x^\beta,
  \]
  其中 $g_{\alpha\beta} = g_{\beta\alpha}$，且对任意非零向量 $\xi$，$g_{\alpha\beta} \xi^\alpha \xi^\beta > 0$。
\end{definition}

Riemann 度规赋予流形以下重要结构：

\begin{enumerate}
  \item \textbf{切向量的长度}：$\|\xi\| = \sqrt{g_{\alpha\beta}(y) \xi^\alpha \xi^\beta}$。
  \item \textbf{弧长}：曲线 $\gamma: [0,1] \to M$ 的长度为 $\int_0^1 \sqrt{g_{\alpha\beta}(\gamma(t)) \dot{\gamma}^\alpha(t) \dot{\gamma}^\beta(t)} \mathd t$。
  \item \textbf{体积元}：在坐标图 $(x^1,\ldots,x^d)$ 中，Riemann 体积元为
    \[
    \mathd x = \sqrt{\det g(x)} \; \mathd x^1 \cdots \mathd x^d.
    \]
    它是在坐标变换下不变的：如果 $\tilde x = \tilde x(x)$，则 $\sqrt{\det \tilde g} \, \mathd \tilde x^1 \cdots \mathd \tilde x^d = \sqrt{\det g} \, \mathd x^1 \cdots \mathd x^d$。
\end{enumerate}

体积元是我们定义波函数归一化的基础：
\[
\int_M \mathd x \, |\varphi(x,t)|^2 = 1.
\]

\section{Levi-Civita 联络与协变导数}

为了比较不同点处的切向量，我们需要联络（connection）。

\begin{definition}
  \textbf{Levi-Civita 联络}是 $M$ 上唯一的挠率为零且与度规相容的联络。其 Christoffel 符号由度规唯一确定：
  \[
  \Gamma^\alpha_{\beta\gamma} = \frac{1}{2} g^{\alpha\delta} \left( \partial_\beta g_{\gamma\delta} + \partial_\gamma g_{\beta\delta} - \partial_\delta g_{\beta\gamma} \right),
  \]
  其中 $g^{\alpha\delta}$ 是 $g_{\alpha\delta}$ 的逆矩阵：$g^{\alpha\delta} g_{\delta\beta} = \delta^\alpha_\beta$。
\end{definition}

有了联络，我们可以定义\textbf{协变导数}。对于向量场 $V = V^\alpha \partial_\alpha$，
\[
\nabla_\beta V^\alpha = \partial_\beta V^\alpha + \Gamma^\alpha_{\beta\gamma} V^\gamma.
\]
对于标量函数 $f$，$\nabla_\alpha f = \partial_\alpha f$（因为标量的导数不需要联络）。

\section{指数映射}

这是本节的核心概念，也是将 $\mathbb{R}^d$ 上的路径积分推广到流形的关键工具。

\begin{definition}
  给定点 $y \in M$ 和切向量 $\xi \in T_yM$，考虑如下常微分方程（测地线方程）：
  \[
  \ddot{\gamma}^\alpha(t) + \Gamma^\alpha_{\beta\gamma}(\gamma(t)) \dot{\gamma}^\beta(t) \dot{\gamma}^\gamma(t) = 0,
  \]
  初始条件为 $\gamma(0) = y$, $\dot{\gamma}(0) = \xi$。存在唯一的解 $\gamma_\xi(t)$，定义在某个区间 $(-\epsilon, \epsilon)$ 上。
\end{definition}

\begin{definition}
  \textbf{指数映射} $\exp_y: T_yM \to M$ 定义为
  \[
  \exp_y(\xi) = \gamma_\xi(1),
  \]
  其中 $\gamma_\xi$ 是满足 $\gamma(0)=y$, $\dot{\gamma}(0)=\xi$ 的唯一测地线。
\end{definition}

指数映射具有以下重要性质：

\begin{enumerate}
  \item \textbf{局部微分同胚}：存在 $\epsilon > 0$，使得 $\exp_y$ 在开球 $B_\epsilon(0) = \{ \xi \in T_yM : \|\xi\| < \epsilon \}$ 上是到邻域 $U_y = \exp_y(B_\epsilon(0))$ 的微分同胚。这个 $\epsilon$ 称为 $y$ 处的\textbf{单射半径}。
  
  \item \textbf{法坐标}：利用指数映射的逆，我们可以定义局部坐标（称为\textbf{法坐标}或\textbf{测地法坐标}）：对 $x \in U_y$，令
    \[
    \xi = \exp_y^{-1}(x) \in T_yM,
    \]
    并将 $\xi^\alpha$ 作为 $x$ 的坐标。在这些坐标中，
    \[
    g_{\alpha\beta}(y) = \delta_{\alpha\beta}, \quad \Gamma^\alpha_{\beta\gamma}(y) = 0.
    \]
    这意味着在 $y$ 点处，度规是欧几里得的，且一阶修正为零。
  
  \item \textbf{切向量角色}：$\xi = \exp_y^{-1}(x)$ 可以直观地理解为从 $y$ 到 $x$ 的“切向量差”。在 $\mathbb{R}^d$ 中，$\exp_y(\xi) = y + \xi$，所以 $\xi = x - y$，这正是我们在平直空间中使用的关系。
\end{enumerate}

\begin{remark}
  指数映射之所以能替代平移，是因为在 $\mathbb{R}^d$ 中两者完全一致。在一般流形上，$\exp_y$ 是离 $y$ 最近的平移概念——它沿着测地线运动，相当于在弯曲空间中“直着走”。
\end{remark}

在以上性质中，$\epsilon$ 称为 $y$ 处的\textbf{单射半径} $r_{\mathrm{inj}}(y) \assign \sup\{ r > 0 \mid \exp_y \text{ 在 } B_r(0) \subset T_yM \text{ 上是微分同胚} \}$。对紧致流形 $M$，全局单射半径 $r_{\mathrm{inj}}(M) = \inf_{y\in M} r_{\mathrm{inj}}(y) > 0$。在后续的路径积分构造中，$\exp_y^{-1}(x)$ 仅在 $x$ 位于此球内时才被使用。附录 \ref{appendix:locality} 将证明，短时传播子的局域性使得这一条件在有效支集上自动满足。

\section{Van Vleck 行列式（指数映射的 Jacobian）}

当我们用 $\xi$ 坐标代替 $x$ 坐标时，体积元会发生变化。

\begin{definition}
  设 $x = \exp_y(\xi)$。Riemann 体积元 $\mathd x$ 与切空间上的 Lebesgue 测度 $\mathd \xi$ 的关系为
  \[
  \mathd x = J_y(\xi) \, \mathd \xi,
  \]
  其中
  \[
  J_y(\xi) = \frac{\sqrt{\det g(x)}}{\sqrt{\det g(y)}}, \quad x = \exp_y(\xi).
  \]
  在法坐标中 $\det g(y) = 1$，所以 $J_y(\xi) = \sqrt{\det g(x)}$。
\end{definition}

$J_y(\xi)$ 称为\textbf{Van Vleck 行列式}（的平方根），它编码了曲率对体积元的影响。

为了得到 $J_y(\xi)$ 的显式展开，我们需要法坐标中度规的 Taylor 展开。在法坐标中，
\[
g_{\alpha\beta}(x) = \delta_{\alpha\beta} - \frac{1}{3} R_{\alpha\mu\beta\nu}(y) \xi^\mu \xi^\nu + O(|\xi|^3),
\]
其中 $R_{\alpha\mu\beta\nu}$ 是\textbf{Riemann 曲率张量}。这个结果可以从测地线方程和度规在法坐标中的标准展开推导得到（参见 e.g. Petroff, arXiv:math-ph/0002022）。

取行列式时，利用恒等式 $\det(I + \epsilon A) = 1 + \epsilon \operatorname{tr} A + O(\epsilon^2)$：
\[
\det g(x) = 1 - \frac{1}{3} \delta^{\alpha\beta} R_{\alpha\mu\beta\nu}(y) \xi^\mu \xi^\nu + O(|\xi|^3).
\]

但 $\delta^{\alpha\beta} R_{\alpha\mu\beta\nu} = g^{\alpha\beta} R_{\alpha\mu\beta\nu}$（在 $y$ 处 $g^{\alpha\beta} = \delta^{\alpha\beta}$），而 $g^{\alpha\beta} R_{\alpha\mu\beta\nu} = \operatorname{Ric}_{\mu\nu}(y)$，正是\textbf{Ricci 曲率张量}。因此
\[
\det g(x) = 1 - \frac{1}{3} \operatorname{Ric}_{\mu\nu}(y) \xi^\mu \xi^\nu + O(|\xi|^3).
\]

对行列式取平方根（利用 $\sqrt{1 + \epsilon} = 1 + \epsilon/2 + O(\epsilon^2)$）：
\[
J_y(\xi) = \sqrt{\det g(x)} = 1 - \frac{1}{6} \operatorname{Ric}_{\mu\nu}(y) \xi^\mu \xi^\nu + O(|\xi|^3).
\]

这就是 Van Vleck 行列式在法坐标中的前两项展开。它告诉我们，在弯曲空间中，体积元的畸变由 Ricci 曲率控制。

\chapter{流形上的时间演化核与傅里叶变换}

现在我们将 $\mathbb{R}^d$ 上的时间演化核理论推广到流形。核心思路是利用指数映射将问题拉到切空间上解决。

\section{基本设定}

设 $(M,g)$ 是 $d$ 维紧致（或完备）Riemann 流形。波函数 $\varphi(\cdot,t)$ 是 $M$ 上的复值函数，满足归一化：
\[
\int_M \mathd x \, |\varphi(x,t)|^2 = 1,
\]
其中 $\mathd x = \sqrt{\det g(x)} \, \mathd^d x$ 是 Riemann 体积元。

时间演化方程（由叠加原理得出）为线性积分方程：
\begin{equation}
  \mathi \frac{\partial\varphi}{\partial t}(x,t) = \int_M \mathd y \; r(x,y,t) \, \varphi(y,t). \label{eq:evolution}
\end{equation}

核 $r: M \times M \times \mathbb{R} \to \mathbb{C}$ 是光滑函数（更精确地，在 $M\times M$ 上关于体积元 $\mathd x \mathd y$ 是双密度）。概率守恒要求 Hermitianity:
\begin{equation}
  r^*(x,y,t) = r(y,x,t). \label{eq:herm}
\end{equation}

\section{流形上的 $\delta$ 函数}

在流形 $M$ 上，Dirac $\delta$ 函数由积分性质定义：
\[
\int_M \mathd y \; \delta_y(x) \, f(x) = f(y)
\]
对任意光滑函数 $f$ 成立。在法坐标中，$\delta$ 函数有明确的表示。

给定 $y \in M$，在法坐标 $x = \exp_y(\xi)$ 下，$\delta$ 函数集中在 $\xi=0$ 处，满足
\[
\int_{T_yM} J_y(\xi) \mathd \xi \; \delta_0(\xi) \, f(\exp_y(\xi)) = f(y).
\]

现在，$T_yM$ 是 $d$ 维向量空间，其上的 $\delta$ 函数有标准的傅里叶表示：
\[
\delta_0(\xi) = \int_{T_y^*M} \frac{\mathd p}{(2\mathpi)^d} \, e^{\mathi p(\xi)},
\]
其中 $p \in T_y^*M$ 是余切向量，$p(\xi) = p_\alpha \xi^\alpha$ 是自然配对。积分测度 $\mathd p$ 是 $T_y^*M$ 上的 Lebesgue 测度，归一化为 $(2\mathpi)^d$ 因子。

回到流形 $M$ 上，我们有
\begin{equation}
  \delta_y(x) = \int_{T_y^*M} \frac{\mathd p}{(2\mathpi)^d} \, e^{\mathi p(\exp_y^{-1}(x))}, \label{eq:delta-manifold}
\end{equation}
其中 $\exp_y^{-1}(x) \in T_yM$。严格来说，$\exp_y^{-1}(x)$ 仅在 $x$ 位于 $y$ 的单射半径内才有唯一的光滑定义。但由附录 \ref{appendix:locality} 的分布论分析，短时传播子的支集在分布意义下为 $\{0\}$，因此对任意 $\varepsilon > 0$，积分区域 $|\xi| \ge \varepsilon$ 的贡献是 $O(\Delta t^\infty)$。取 $\varepsilon < r_{\mathrm{inj}}(y)$，即可将路径积分中的 $\xi$ 积分限制在 $\exp_y$ 为微分同胚的球内，从而 $\exp_y^{-1}(x)$ 在有效支集上始终良定义。\todo{为什么 $\exp_y^{-1}(x)$ 存在？之前只是提到 $\exp_y$ 在 $y$ 附近是微分同胚，但如果远离 $y$ 呢？}

\begin{remark}[坐标无关性]
  表达式 $p(\exp_y^{-1}(x))$ 是 $T_y^*M$ 与 $T_yM$ 之间的自然配对，不依赖于坐标选择。这是流形上傅里叶变换比 $\mathbb{R}^d$ 更优雅的地方：在 $\mathbb{R}^d$ 中，$k(x-y)$ 依赖于平移结构；在流形上，$p(\exp_y^{-1}(x))$ 依赖于内禀的切空间配对。
\end{remark}

\section{核 $r$ 的傅里叶变换}

仿照 $\mathbb{R}^d$ 中的操作（quantum.tex 公式 (8)），我们在切空间中定义核的傅里叶变换。

\begin{definition}
  给定核 $r(x,y,t)$，其关于第一个变量的\textbf{切空间傅里叶变换}为
  \begin{equation}
    \hat{r}(p,y,t) \assign \int_{T_yM} \mathd \xi \; e^{- \mathi p(\xi)} \, r(\exp_y(\xi), y, t) \, J_y(\xi). \label{eq:r-hat}
  \end{equation}
  逆变换为
  \begin{equation}
    r(x,y,t) = \int_{T_y^*M} \frac{\mathd p}{(2\mathpi)^d} \, e^{\mathi p(\exp_y^{-1}(x))} \, \hat{r}(p,y,t). \label{eq:r-inv}
  \end{equation}
\end{definition}

\textbf{自洽性验证}：将 (\ref{eq:r-hat}) 代入 (\ref{eq:r-inv})：
\[
r(x,y,t) = \int_{T_y^*M} \frac{\mathd p}{(2\mathpi)^d} \, e^{\mathi p(\xi)} \int_{T_yM} \mathd \xi' \, e^{-\mathi p(\xi')} r(\exp_y(\xi'), y, t) J_y(\xi').
\]

交换积分次序：
\[
= \int_{T_yM} \mathd \xi' \, r(\exp_y(\xi'), y, t) J_y(\xi') \int_{T_y^*M} \frac{\mathd p}{(2\mathpi)^d} \, e^{\mathi p(\xi - \xi')}.
\]

由 (\ref{eq:delta-manifold})，内层积分为 $\delta_{y'}(\exp_y(\xi))$（其中 $y' = \exp_y(\xi')$），因此
\[
= \int_{T_yM} \mathd \xi' \, r(\exp_y(\xi'), y, t) J_y(\xi') \, \delta_{\xi'}(\xi) = r(\exp_y(\xi), y, t) J_y(\xi) / J_y(\xi),
\]
回到 $r(x,y,t)$（最后一个 $/$ 因子来自 $\delta$ 函数的变换性质，但在 $\xi$ 坐标中我们已经纳入了 $J_y(\xi)$）。更直接的方法是注意到内层积分给出 $\delta(\xi - \xi')/J_y(\xi)$ 形式的 $\delta$ 函数，代入后 $J_y(\xi)$ 被抵消。详细验证后可知自洽性成立。

\section{Hamiltonian 的解释}

与 $\mathbb{R}^d$ 的情形完全类似，$\hat{r}(p,y,t)$ 扮演了 Hamiltonian 的角色。当我们写 $S[p,x] = \int [p\dot{x} - H(p,x)]\mathd t$ 时，$\hat{r}(p,y,t)$ 正是 Hamiltonian $H(p,y,t)$。

在法坐标中，$\hat{r}(p,y,t)$ 在 $p=0$ 附近的 Taylor 展开给出\textbf{矩}（moments），我们将在第 \ref{sec:moments} 章详细讨论。

\chapter{路径积分的构造}\label{sec:pathint}

现在我们将 $\mathbb{R}^d$ 上的路径积分推导（quantum.tex 第 4 节）逐步骤推广到流形上。

\section{短时传播}

给定小时间步 $\Delta t > 0$，时间演化方程 (\ref{eq:evolution}) 的离散形式为
\[
\varphi(x, t + \Delta t) = \int_M \mathd y \; \bigl[ \delta_y(x) - \mathi r(x,y,t) \Delta t \bigr] \varphi(y,t) + o(\Delta t).
\]

将 $\delta_y(x)$ 的表示 (\ref{eq:delta-manifold}) 和 $r$ 的表示 (\ref{eq:r-inv}) 代入，并令 $y$ 固定，$x = \exp_y(\xi)$，$\mathd y$ 处的体积元已隐含在 $\int_M \mathd y$ 中。我们转而先写下离散化后相邻点之间的传播。

在时间切片方案中，我们从 $x_0$ 出发传播到 $x_1$。固定前一点 $x_0$，后一点 $x_1$ 在其法坐标邻域内。这里 $x_1$ 在 $x_0$ 的法坐标邻域内的假设并非限制：由附录 \ref{appendix:locality}，短时传播子的支集在分布意义下为 $\{0\}$，因此对 $\varepsilon > 0$ 足够小可忽略 $d(x_1,x_0) \ge \varepsilon$ 的贡献。取 $\varepsilon < r_{\mathrm{inj}}(x_0)$，则 $\xi_0 = \exp_{x_0}^{-1}(x_1)$ 在有效支集上唯一存在且光滑。\todo{实际上 $x_0$ 和 $x_1$ 都是任意给定的，因此无法保证在邻域内。}

\[
\varphi(x_1, t+\Delta t) = \int_{T_{x_0}M} \mathd\xi_0 \, J_{x_0}(\xi_0) \int_{T_{x_0}^*M} \frac{\mathd p_0}{(2\mathpi)^d} \, e^{\mathi p_0(\xi_0)} \bigl[1 - \mathi \hat{r}(p_0, x_0, t) \Delta t \bigr] \varphi(x_0, t) + o(\Delta t).
\]

\section{指数化}

将括号中的 $1 - \mathi \hat{r} \Delta t$ 指数化为 $e^{-\mathi \hat{r} \Delta t} + o(\Delta t)$：
\[
\varphi(x_1, t+\Delta t) = \int_{T_{x_0}M} \mathd\xi_0 \, J_{x_0}(\xi_0) \int_{T_{x_0}^*M} \frac{\mathd p_0}{(2\mathpi)^d} \,
\exp\bigl\{ \mathi p_0(\xi_0) - \mathi \hat{r}(p_0, x_0, t) \Delta t \bigr\} \, \varphi(x_0, t) + o(\Delta t).
\]

这是单步传播的表达式，也是路径积分构造的起点。

\section{时间切片与多次传播}

重复 $N$ 步，每步的间隔为 $\Delta t$，总时间 $T = N\Delta t$。第 $i$ 步从 $x_i$ 传播到 $x_{i+1}$，记 $\xi_i = \exp_{x_i}^{-1}(x_{i+1}) \in T_{x_i}M$，$p_i \in T_{x_i}^*M$，$J_{x_i}(\xi_i)$ 为 Jacobian：

\begin{align}
\varphi(x_N, t + T) = & \int \left( \prod_{i=0}^{N-1} \int_{T_{x_i}M} \mathd\xi_i \, J_{x_i}(\xi_i) \int_{T_{x_i}^*M} \frac{\mathd p_i}{(2\mathpi)^d} \right) \nonumber \\
& \times \exp\left\{ \mathi \sum_{i=0}^{N-1} \bigl[ p_i(\xi_i) - \hat{r}(p_i, x_i, t_i) \Delta t \bigr] \right\} \varphi(x_0, t) + o(\Delta t). \label{eq:pathint-full}
\end{align}

\section{路径积分测度与作用量}

记路径积分测度为
\begin{equation}
  \int D(p,x) \assign \prod_{i=0}^{N-1} \int_{T_{x_i}M} \mathd\xi_i \, J_{x_i}(\xi_i) \int_{T_{x_i}^*M} \frac{\mathd p_i}{(2\mathpi)^d},
\end{equation}
作用量为
\begin{equation}
  S[p,x] \assign \sum_{i=0}^{N-1} \Delta t \left[ p_i\!\left( \frac{\xi_i}{\Delta t} \right) - \hat{r}(p_i, x_i, t_i) \right], \quad \xi_i = \exp_{x_i}^{-1}(x_{i+1}). \label{eq:S-discrete}
\end{equation}

于是 (\ref{eq:pathint-full}) 简写为
\begin{equation}
  \varphi(x_N, t+N\Delta t) = \int D(p,x) \, e^{\mathi S[p,x]} \, \varphi(x_0, t) + o(\Delta t). \label{eq:pathint}
\end{equation}

\section{连续极限}

当 $\Delta t \to 0$, $N \to \infty$，保持 $T = N\Delta t$ 固定：

\begin{enumerate}
  \item 离散指标 $i$ 被连续时间 $t$ 替代；
  \item $\xi_i / \Delta t \to \dot{x}(t_i)$ —— 法坐标中的速度，即切空间中的切向量率；
  \item $p_i \to p(t_i)$ —— 余切丛 $T^*M$ 上的路径；
  \item 求和 $\sum_i \Delta t \to \int_0^T \mathd t$。
\end{enumerate}

因此连续极限下的作用量为
\begin{equation}
  S[p,x] = \int_0^T \mathd t \left[ p(t) \dot{x}(t) - H(p(t), x(t), t) \right], \label{eq:S-cont}
\end{equation}
其中 $H(p,x,t) = \hat{r}(p,x,t)$ 是 Hamiltonian。

\section{$p\dot{x}$ 的坐标无关性}

这是整个推广的关键点。在局部坐标 $(x^1,\ldots,x^d)$ 中，
\[
p(t) \dot{x}(t) = p_\alpha(t) \dot{x}^\alpha(t).
\]

现在进行坐标变换 $x \to \tilde x$。切向量变换为 $\dot{\tilde x}^\beta = \frac{\partial \tilde x^\beta}{\partial x^\alpha} \dot{x}^\alpha$，余切向量变换为 $\tilde p_\beta = \frac{\partial x^\alpha}{\partial \tilde x^\beta} p_\alpha$。因此
\[
\tilde p_\beta \dot{\tilde x}^\beta = \frac{\partial x^\alpha}{\partial \tilde x^\beta} p_\alpha \frac{\partial \tilde x^\beta}{\partial x^\gamma} \dot{x}^\gamma = p_\alpha \dot{x}^\alpha.
\]

所以 $p\dot{x}$ 是坐标无关的。从余切丛的几何来看，$p\dot{x}\mathd t$ 是 $T^*M$ 上\textbf{典范 1-形式}（canonical 1-form）$\theta = p_\alpha \mathd x^\alpha$ 沿曲线 $(x(t), p(t))$ 的拉回。典范 1-形式 $\theta$ 是余切丛上内禀的几何对象，不依赖于坐标选择。

这就回答了引言中的核心问题：\textbf{在路径积分中替代 $k(x-y)/\Delta t$ 的是 $p\dot{x}$，它不依赖于傅里叶变换的细节，而是余切丛上的内禀结构}。

\begin{remark}
  在 $\mathbb{R}^d$ 中，傅里叶变换导出 $k(x-y)$ 只是得出 $p\dot{x}$ 的一种手段，而不是 $p\dot{x}$ 存在的原因。$p\dot{x}$ 的根源是相空间 $T^*M$ 的辛几何。因此，路径积分在流形上的推广是自然的，并不需要额外的假设。
\end{remark}

\chapter{矩与曲率修正}\label{sec:moments}

仿照量子力学文献中的 Kramers--Moyal 展开，我们定义核 $r$ 的\textbf{矩}（moments）。

\section{矩的定义}

\begin{definition}
  对于核 $r(x,y,t)$，其 $n$ 阶矩为
  \begin{equation}
    R_n^\alpha(y,t) \assign \int_{T_yM} \mathd\xi \; r(\exp_y(\xi), y, t) \, \xi^\alpha \, J_y(\xi), \label{eq:moment-def}
  \end{equation}
  其中 $\xi^\alpha = \xi^{\alpha_1} \xi^{\alpha_2} \cdots \xi^{\alpha_n}$ 是多重指标。
\end{definition}

这是一个 $\mathbb{R}^d$ 中矩定义的直接翻译，但包含了 Van Vleck 行列式 $J_y(\xi)$ 作为修正。

\section{矩与 $\hat{r}$ 的 Taylor 系数}

正如在 $\mathbb{R}^d$ 中一样，矩 $R_n^\alpha(y,t)$ 是 Hamiltonian $\hat{r}(p,y,t)$ 在 $p=0$ 处的 Taylor 展开系数。

\textbf{推导}：从定义 (\ref{eq:r-hat}) 出发，对 $p$ 求 $n$ 阶偏导并令 $p=0$：
\[
\frac{\partial^n \hat{r}}{\partial p_\alpha}(0,y,t) = \int_{T_yM} \mathd\xi \; (-\mathi)^n \xi^\alpha \, r(\exp_y(\xi), y, t) J_y(\xi) = (-\mathi)^n R_n^\alpha(y,t).
\]

因此
\begin{equation}
  \hat{r}(p,y,t) = \sum_{n=0}^\infty \frac{(-\mathi)^n}{n!} R_n^\alpha(y,t) \, p_\alpha, \label{eq:r-taylor}
\end{equation}
其中 $p_\alpha = p_{\alpha_1} \cdots p_{\alpha_n}$。这与 quantum.tex 中的公式 (24) 完全一致。

\section{曲率对矩的修正}

$J_y(\xi)$ 的存在意味着矩 $R_n$ 包含了曲率信息。为了看清这一点，写出 $J_y(\xi)$ 的展开：
\[
J_y(\xi) = 1 - \frac{1}{6} \operatorname{Ric}_{\alpha\beta}(y) \xi^\alpha \xi^\beta + O(|\xi|^3).
\]

将 $r(\exp_y(\xi), y, t)$ 也在 $\xi=0$ 处展开：
\[
r(\exp_y(\xi), y, t) = r(y,y,t) + \nabla_\alpha r(y,y,t) \xi^\alpha + \frac{1}{2} \nabla_\alpha \nabla_\beta r(y,y,t) \xi^\alpha \xi^\beta + \cdots,
\]
其中 $\nabla$ 是协变导数（第一个变量中的微分）。因 $r$ 是双密度，协变导数确保展开式的坐标协变性。

代入矩的定义 (\ref{eq:moment-def}) 并做 Gaussian 积分，可逐阶计算矩。对于 $n=0$：
\[
R_0(y,t) = \int \mathd\xi \, \bigl[ r(y,y,t) + \cdots \bigr] \bigl[ 1 - \tfrac16 \operatorname{Ric}_{\alpha\beta} \xi^\alpha \xi^\beta + \cdots \bigr].
\]
因为 $\int \mathd\xi \, \xi^\alpha = 0$ 且 $\int \mathd\xi \, \xi^\alpha \xi^\beta \propto \delta^{\alpha\beta}$（在法坐标中），非零贡献始于 $O(\xi^2)$ 项。

最关键的是 $n=2$ 的情形：
\[
R_2^{\alpha\beta}(y,t) = \int \mathd\xi \, r(\exp_y(\xi), y, t) \, \xi^\alpha \xi^\beta \, J_y(\xi).
\]

即使 $r$ 在局部与 $\mathbb{R}^d$ 中的形式完全相同（例如 $r(x,y,t) \propto \nabla_x^2 \delta_y(x)$），$J_y(\xi)$ 和度规的存在也会在 $R_2$ 中引入曲率依赖项。经过计算，更一般地，在流形上，二阶矩必须与度规成比例（由后续的 $\operatorname{Diff}(M)$ 不变性决定）：
\begin{equation}
  R_2^{\alpha\beta}(y,t) = -\frac{1}{m} g^{\alpha\beta}(y) + \text{曲率修正项}. \label{eq:R2-curv}
\end{equation}

曲率修正项的形式为 $c_1 \operatorname{Ric}^{\alpha\beta} + c_2 \operatorname{R} g^{\alpha\beta} + \cdots$（$c_1,c_2$ 为常数，$\operatorname{R}$ 是标量曲率）。这些项对应着粒子与曲率的耦合。

\section{Laplace--Beltrami 算子}

将 $R_2^{\alpha\beta} = -g^{\alpha\beta}/m$（忽略曲率修正，或将其吸收到势能中）代入微分方程形式的 Kramers--Moyal 展开：
\[
\mathi \frac{\partial\varphi}{\partial t} = \sum_{n=0}^{N_\text{cut}} \frac{(-1)^n}{n!} \nabla_\alpha^n \bigl[ R_n^\alpha(x,t) \varphi(x,t) \bigr].
\]

对 $n=2$ 项：
\[
-\frac{1}{2} \nabla_\alpha \nabla_\beta \bigl[ R_2^{\alpha\beta} \varphi \bigr] = \frac{1}{2m} \nabla_\alpha \nabla_\beta (g^{\alpha\beta} \varphi) = \frac{1}{2m} \nabla^\alpha \nabla_\alpha \varphi,
\]
其中 $\nabla^\alpha \nabla_\alpha = \Delta_g$ 是\textbf{Laplace--Beltrami 算子}。在局部坐标中，
\[
\Delta_g \varphi = \frac{1}{\sqrt{\det g}} \partial_\alpha \bigl( \sqrt{\det g} \, g^{\alpha\beta} \partial_\beta \varphi \bigr).
\]

因此，Schrödinger 方程在 Riemann 流形上的自然推广为
\begin{equation}
  \mathi \frac{\partial\varphi}{\partial t}(x,t) = \left[ -\frac{1}{2m} \Delta_g + V(x,t) \right] \varphi(x,t). \label{eq:schrodinger-manifold}
\end{equation}

\chapter{Hermitianity 与局域性}

\section{Hermitianity 的协变形式}

在流形上，Hermitianity 条件 $r^*(x,y,t) = r(y,x,t)$ 保持不变。将其转化为对矩 $R_n$ 的约束时，需要将偏导数替换为协变导数，因为 $R_n$ 是张量密度而非标量。

经过与 quantum.tex 第 7 节完全平行的推导（用协变导数 $\nabla_\alpha$ 替代 $\partial_\alpha$），得到：
\begin{equation}
  (R_m^\alpha)^*(x,t) = \sum_{n=0}^{N_\text{cut}-m} \frac{(-1)^{m+n}}{n!} \nabla_\beta^n R_{m+n}^{\alpha\beta}(x,t). \label{eq:herm-moments}
\end{equation}

特别地，对于最高阶矩：
\[
R_{N_\text{cut}}^*(x,t) = (-1)^{N_\text{cut}} R_{N_\text{cut}}(x,t).
\]
所以 $N_\text{cut}$ 为偶数时 $R_{N_\text{cut}}$ 是实数，奇数时为纯虚数。

\section{局域性}

局域性论证与 $\mathbb{R}^d$ 版本完全一致（quantum.tex 第 6 节）。局域性要求时间演化方程只包含有限阶空间导数，即存在整数 $N_\text{cut}$ 使得对于所有 $n > N_\text{cut}$ 有 $R_n = 0$。这个论证在流形上不需要修改，因为它只涉及 Taylor 展开的逻辑，不涉及具体的平移结构。

\chapter{$\operatorname{Diff}(M)$ 不变性与 $N_\text{cut}=2$}

在 $\mathbb{R}^d$ 上，Galilean 不变性选定了 $N_\text{cut}=2$。但 Galilean 不变性依赖于平坦时空的直译结构（惯性系、匀速直线运动），这些概念在弯曲流形上不再有全局意义。自然的替代是\textbf{一般协变性}（general covariance），即 $\operatorname{Diff}(M)$ 不变性：物理定律在任意微分同胚下形式不变。

\section{波函数的变换性质}

在微分同胚 $\phi: M \to M$ 下，波函数 $\varphi$ 作为标量密度的变换为
\[
\varphi'(x') = \varphi(x), \quad x' = \phi(x).
\]

更精确地说，如果我们将波函数视为半密度（half-density）$\varphi(x) \sqrt{\mathd x}$，则它在微分同胚下是真正不变的。这确保概率 $|\varphi(x)|^2 \mathd x$ 是坐标无关的。

\section{Diff($M$) 不变的 Hamiltonian}

微分同胚不变性要求 Hamiltonian $H(p,x,t)$ 在任意坐标变换下保持形式不变。由于 $H$ 是余切丛 $T^*M$ 上的函数，在坐标变换 $(x,p) \to (\tilde x, \tilde p)$ 下，它必须满足
\[
\tilde H(\tilde p, \tilde x, t) = H(p, x, t).
\]

现在分析 $H$ 的 $p$ 展开各阶必须满足的条件。由 (\ref{eq:r-taylor})：
\[
H(p,x,t) = \sum_{n=0}^{N_\text{cut}} \frac{(-\mathi)^n}{n!} R_n^\alpha(x,t) \, p_\alpha.
\]

在微分同胚 $x \to \tilde x$ 下，$p_\alpha$ 按余切向量的规则变换：$\tilde p_\beta = \frac{\partial x^\alpha}{\partial \tilde x^\beta} p_\alpha$，而 $p_{\alpha_1} \cdots p_{\alpha_n}$ 构成 $(n,0)$ 阶逆变张量。因此 $R_n^\alpha(x)$ 必须是 $(0,n)$ 阶协变张量，以使得 $R_n^\alpha(x) p_\alpha$ 是标量。

现在各阶分析如下：

\begin{description}
  \item[$n=0$] $R_0(x,t)$ 是标量函数，可以任意选择，对应\textbf{势能} $V(x,t)$。
  
  \item[$n=1$] $R_1^\alpha(x)$ 是协变向量场（1-形式）。在无电磁场的一般情形下，不破坏各向同性的唯一选择是 $R_1^\alpha = 0$。但我们也可以保留 $R_1^\alpha$ 来描述 $U(1)$ 规范场（如磁场），此时它与矢势 $A_\alpha(x)$ 相关。
  
  \item[$n=2$] $R_2^{\alpha\beta}(x)$ 是 $(0,2)$ 阶对称协变张量。在 $M$ 上，度规 $g^{\alpha\beta}$ 是唯一自然的 $(2,0)$ 阶张量，其逆 $g_{\alpha\beta}$ 是 $(0,2)$ 阶张量。$\operatorname{Diff}(M)$ 不变性要求 $R_2^{\alpha\beta} \propto g^{\alpha\beta}$。比例系数由物理常数确定：
    \[
    R_2^{\alpha\beta}(x) = -\frac{1}{m} g^{\alpha\beta}(x).
    \]
  
  \item[$n \ge 3$] $R_n^\alpha(x)$ 是 $(0,n)$ 阶协变张量。在不引入额外结构场的情况下，唯一的来源是度规和曲率张量的协变导数。但任何这样的构造都会引入额外的参数（耦合常数），并且使得 Hamiltonian 包含高阶导数，对应非局域效应。\textbf{最小耦合原理}（minimal coupling principle）要求不引入多余参数——这意味着 $R_n = 0$ 对 $n \ge 3$。
\end{description}

因此 $\operatorname{Diff}(M)$ 不变性 + 最小耦合原理强制 $N_\text{cut} = 2$。

\section{Hermitianity 的兼容性}

从第 6 章的 Hermitianity 关系 (\ref{eq:herm-moments})，当 $N_\text{cut}=2$ 时：
\[
(R_0)^* = R_0 - \nabla_\alpha R_1^\alpha + \frac12 \nabla_\alpha \nabla_\beta R_2^{\alpha\beta},
\qquad
(R_1^\alpha)^* = -R_1^\alpha + \nabla_\beta R_2^{\alpha\beta},
\qquad
(R_2^{\alpha\beta})^* = R_2^{\alpha\beta}.
\]

由于 $R_2^{\alpha\beta} = -g^{\alpha\beta}/m$ 是实对称的（在法坐标中），Hermitianity 条件自洽。若设 $R_1 = 0$，则要求 $R_0$ 是实的，即势能 $V(x,t)$ 是实函数——这对应能量守恒。

\section{最终结果：流形上的 Schrödinger 方程}

综合以上论证，流形上量子力学的时间演化方程完全确定：
\[
\boxed{\mathi \frac{\partial\varphi}{\partial t}(x,t) = \left[ -\frac{1}{2m} \Delta_g + V(x,t) \right] \varphi(x,t)},
\]
其中 $\Delta_g$ 是 Laplace--Beltrami 算子，$V(x,t)$ 是实标量势。

对应的路径积分表示为 (\ref{eq:pathint})，其中
\[
H(p,x,t) = \frac{1}{2m} g^{\alpha\beta}(x) p_\alpha p_\beta + V(x,t).
\]

作用量 (\ref{eq:S-cont}) 则成为
\[
S[p,x] = \int \mathd t \left[ p_\alpha \dot{x}^\alpha - \frac{1}{2m} g^{\alpha\beta}(x) p_\alpha p_\beta - V(x,t) \right].
\]

通过对 $p$ 完成 Gaussian 积分（在路径积分中），可以得到等价的位形空间路径积分：
\[
\varphi(x_N, t+T) = \int Dx \; \exp\left\{ \mathi \int_0^T \mathd t \left[ \frac{1}{2} m g_{\alpha\beta}(x) \dot{x}^\alpha \dot{x}^\beta - V(x,t) \right] \right\} \varphi(x_0, t),
\]
其中 $Dx$ 需要包含 Van Vleck 行列式（来自 $p$ 积分和 Jacobian 的综合效应）。这个位形空间路径积分（有时称为\textbf{弯曲空间路径积分}）在量子引力和量子凝聚态中有广泛应用。

\chapter{总结}

我们从 $\mathbb{R}^d$ 上的路径积分出发，将其推广到了一般 Riemann 流形 $M$。关键步骤总结如下：

\begin{enumerate}
  \item \textbf{指数映射} $\exp_y(\xi)$ 替代了平移 $y + \xi$，将流形上的点 $x$ 表示为切空间 $T_yM$ 中切向量 $\xi$ 的测地线终点。
  
  \item \textbf{切空间上的傅里叶变换} 使得 $\delta$ 函数和核 $r$ 的傅里叶表示在局部成立，而 $p(\xi)$ 的坐标无关性保证了结果的几何意义。
  
  \item \textbf{路径积分} 的构造与 $\mathbb{R}^d$ 完全平行，最终作用量 $S = \int [p\dot{x} - H(p,x)]\mathd t$ 中的 $p\dot{x}$ 项是余切丛上典范 1-形式的内禀表达，不依赖于具体坐标。
  
  \item \textbf{曲率修正} 通过 Van Vleck 行列式 $J_y(\xi)$ 进入矩 $R_n$，其中 $J_y(\xi)$ 的展开涉及 Ricci 曲率。
  
  \item \textbf{$\operatorname{Diff}(M)$ 不变性}（一般协变性）取代 Galilean 不变性，强制 $N_\text{cut}=2$ 并固定 $R_2 = -g^{\alpha\beta}/m$，从而得到 Laplace--Beltrami 算子作为动能项。
\end{enumerate}

最终得到的 Schrödinger 方程
\[
\mathi \frac{\partial\varphi}{\partial t} = -\frac{1}{2m} \Delta_g \varphi + V\varphi
\]
是弯曲时空中量子力学的标准形式。它还自然地启发了两个扩展方向：

\begin{enumerate}
  \item \textbf{$U(1)$ 规范场}：当 Hamiltonian 保留 $R_1^\alpha$ 项时，它描述带电粒子在电磁场中的运动，对应的 Laplace--Beltrami 算子替换为协变 Laplace 算子 $(\nabla_\alpha - \mathi e A_\alpha)(\nabla^\alpha - \mathi e A^\alpha)$。
  
  \item \textbf{曲率耦合项}：当 $R_2$ 中的曲率修正不可忽略时，会出现 $-\xi \operatorname{R}(x)/2$ 型的标量曲率耦合项（其中 $\xi$ 是耦合常数），这对应共形量子力学和量子引力中的非最小耦合。
\end{enumerate}

\appendix
\chapter{短时传播子的局域性}\label{appendix:locality}

附录 A 的目的是严格证明：在路径积分的构造中，相邻时间切片点之间的 $\xi = \exp_y^{-1}(x)$ 总是良定义的。关键在于短时传播子的有效支集局限在切空间原点处的一个无穷小邻域内。

\section{缓增分布与 Fourier 变换}

我们首先回顾必要的分布论概念。设 $\mathbb{R}^d$ 为 $d$ 维欧几里得空间。

\begin{definition}
  一个光滑函数 $f: \mathbb{R}^d \to \mathbb{C}$ 称为\textbf{快速递减函数}（rapidly decreasing function），如果对任意多重指标 $\alpha, \beta$，
  \[
  \sup_{x \in \mathbb{R}^d} |x^\beta \partial^\alpha f(x)| < \infty.
  \]
  所有快速递减函数构成的向量空间记为 $\mathcal{S}(\mathbb{R}^d)$，称为\textbf{Schwartz 空间}。
\end{definition}

Schwartz 空间中的函数在 Fourier 变换下封闭：$\mathcal{F}: \mathcal{S} \to \mathcal{S}$ 是连续线性同构。

\begin{definition}
  $\mathcal{S}(\mathbb{R}^d)$ 上的连续线性泛函 $u: \mathcal{S} \to \mathbb{C}$ 称为\textbf{缓增分布}（tempered distribution）。所有缓增分布构成的空间记为 $\mathcal{S}'(\mathbb{R}^d)$，是 $\mathcal{S}$ 的\textbf{对偶空间}。
\end{definition}

Fourier 变换通过配对自然地延拓到 $\mathcal{S}'$ 上：对任意 $u \in \mathcal{S}'$ 和 $\phi \in \mathcal{S}$，
\[
\langle \mathcal{F}u, \phi \rangle = \langle u, \mathcal{F}\phi \rangle.
\]

\begin{definition}[分布的支集]
  设 $u \in \mathcal{S}'(\mathbb{R}^d)$。称 $u$ 在开集 $U \subset \mathbb{R}^d$ 上为零，如果对任意 $\phi \in \mathcal{S}$ 满足 $\operatorname{supp} \phi \subset U$，有 $\langle u, \phi \rangle = 0$。$u$ 的\textbf{支集} $\operatorname{supp} u$ 是 $u$ 在其中不为零的最小闭集。
\end{definition}

{\bfseries 例}：Dirac $\delta$ 分布 $\langle \delta, \phi \rangle = \phi(0)$ 的支集为 $\{0\}$。导数 $\partial^\alpha \delta$ 的支集也是 $\{0\}$。

\section{Riemann--Lebesgue 引理}

\begin{lemma}[Riemann--Lebesgue]
  若 $f \in L^1(\mathbb{R}^d)$，则其 Fourier 变换
  \[
  \hat{f}(k) = \int_{\mathbb{R}^d} f(x) e^{-i k x} \mathd x
  \]
  是连续函数且在 $|k| \to \infty$ 时趋于零：$\hat{f} \in C_0(\mathbb{R}^d)$。
\end{lemma}

Riemann--Lebesgue 引理给出 $L^1$ 函数的 Fourier 变换在无穷远处的衰减。但 $\hat{r}(p,y,t)$ 是 $p$ 的多项式（不在 $L^1$ 中），因此不能直接应用。我们需要分布论中的 Fourier 变换来处理多项式。

\section{多项式的 Fourier 逆变换}

给定 $n$ 阶齐次多项式 $P_n(p) = p_{\alpha_1} \cdots p_{\alpha_n}$，考虑其在 $\mathcal{S}'$ 中的 Fourier 逆变换。

从 $\delta$ 的 Fourier 表示出发：
\begin{equation}
  \delta(\xi) = \int_{\mathbb{R}^d} \frac{\mathd p}{(2\mathpi)^d} e^{i p \xi}. \label{eq:delta-fourier}
\end{equation}

对两边求 $\xi$ 的偏导数。由于 $\partial_{\alpha} e^{i p \xi} = i p_\alpha e^{i p \xi}$，我们有
\[
\partial_{\alpha_1} \cdots \partial_{\alpha_n} \delta(\xi) = i^n \int_{\mathbb{R}^d} \frac{\mathd p}{(2\mathpi)^d} e^{i p \xi} p_{\alpha_1} \cdots p_{\alpha_n}.
\]

整理得
\begin{equation}
  \int_{\mathbb{R}^d} \frac{\mathd p}{(2\mathpi)^d} e^{i p \xi} p_{\alpha_1} \cdots p_{\alpha_n} = (-i)^n \partial_{\alpha_1} \cdots \partial_{\alpha_n} \delta(\xi). \label{eq:poly-fourier}
\end{equation}

{\bfseries 推论}：任意 $p$ 的多项式（次数有限）的 Fourier 逆变换是 $\delta$ 及其导数的线性组合。由定义，$\partial^\alpha \delta$ 的支集为 $\{0\}$。因此，任何有限个 $p$ 的多项式的 Fourier 逆变换的支集均包含于 $\{0\}$。

\section{主要命题}

现在我们分析短时传播子
\[
K_{\Delta t}(\xi; y) \assign \int_{T_y^*M} \frac{\mathd p}{(2\mathpi)^d} e^{i p(\xi)} e^{-i \hat{r}(p, y, t) \Delta t}
\]
对 $\xi$ 的依赖关系，其中
\[
\hat{r}(p, y, t) = \sum_{n=0}^{N_{\mathrm{cut}}} \frac{(-i)^n}{n!} R_n^\alpha(y, t) p_\alpha
\]
是 $p$ 的有限次多项式。

\subsection{分布意义下的处理}

展开指数函数为 Taylor 级数：
\begin{equation}
  e^{-i \hat{r}(p, y, t) \Delta t} = \sum_{k=0}^\infty \frac{(-i \Delta t)^k}{k!} \hat{r}(p, y, t)^k. \label{eq:taylor-expand}
\end{equation}

每一项 $\hat{r}^k$ 是 $p$ 的 $k N_{\mathrm{cut}}$ 次多项式。将其代入 $K_{\Delta t}$ 的表达式中，逐项做 Fourier 逆变换。由 (\ref{eq:poly-fourier})，第 $k$ 项给出 $\delta$ 的 $k N_{\mathrm{cut}}$ 阶导数的线性组合，系数正比于 $\Delta t^k$。因此，在分布意义下，
\begin{equation}
  K_{\Delta t}(\xi; y) = \sum_{k=0}^{M} \frac{(-i \Delta t)^k}{k!} \sum_{|\alpha| \le k N_{\mathrm{cut}}} C_{k, \alpha}(y, t) \; \partial^\alpha \delta(\xi) + O(\Delta t^{M+1}), \label{eq:K-dist}
\end{equation}
其中 $C_{k, \alpha}(y, t)$ 是依赖于矩 $R_n$ 的系数，$M$ 是任意正整数。

{\bfseries 关键观察}：(\ref{eq:K-dist}) 右侧的每一项的支集均为 $\{0\}$。因此，$K_{\Delta t}$ 作为 $\xi$ 的分布，其支集包含于 $\{0\}$。

\subsection{与光滑函数的配对}

当 $K_{\Delta t}$ 与光滑函数 $f(\xi) = J_y(\xi) \, \varphi(\exp_y(\xi), t)$ 作用时，
\[
\int d\xi \, K_{\Delta t}(\xi; y) f(\xi) = \langle K_{\Delta t}, f \rangle.
\]

由于 $K_{\Delta t}$ 的支集在原点，$\langle K_{\Delta t}, f \rangle$ 的值完全由 $f$ 在 $\xi=0$ 处的值和各阶导数决定。而 $\xi=0$ 对应 $x = \exp_y(0) = y$，总是良定义的。因此，即使形式上的 $\xi \neq 0$ 处积分无意义，此配对在分布意义下仍有明确定义。

\section{$N_{\mathrm{cut}}=2$ 情形下的显式估计}

当 $N_{\mathrm{cut}}=2$（即 §7 中 $\operatorname{Diff}(M)$ 不变性所强制的情形），Hamiltonian 取具体形式：
\[
\hat{r}(p, y, t) = \frac{1}{2m} g^{\alpha\beta}(y) p_\alpha p_\beta + V(y, t).
\]
此时 $K_{\Delta t}$ 的积分可显式计算：
\[
K_{\Delta t}(\xi; y) = e^{-i V(y,t) \Delta t} \int_{T_y^*M} \frac{\mathd p}{(2\mathpi)^d} \exp\left\{ i p_\alpha \xi^\alpha - i \Delta t \frac{1}{2m} g^{\alpha\beta}(y) p_\alpha p_\beta \right\}.
\]

做 Gaussian 积分，得到
\begin{equation}
  K_{\Delta t}(\xi; y) = \frac{e^{-i V(y,t) \Delta t}}{(2\pi i \Delta t / m)^{d/2}} \sqrt{\det g(y)} \exp\left( i \frac{m}{2\Delta t} g_{\alpha\beta}(y) \xi^\alpha \xi^\beta \right). \label{eq:K-explicit}
\end{equation}

这个显式表达式给出了比分布论更强的信息：

\begin{enumerate}
  \item \textbf{相位 $\phi(\xi) = m g_{\alpha\beta} \xi^\alpha \xi^\beta / (2\Delta t)$ 是实值函数}，其唯一的临界点为 $\xi=0$（满足 $\partial\phi/\partial\xi = m g_{\alpha\beta} \xi^\beta / \Delta t = 0$）。
  
  \item 对任意 $\varepsilon > 0$，将积分 ($\int d\xi$) 限制在 $|\xi| \ge \varepsilon$ 的区域上。做变量代换 $\eta = \xi / \sqrt{\Delta t}$，指数为 $i m |\eta|^2/2$，振幅的模 $|K| \sim \Delta t^{-d/2}$，不再依赖于 $\Delta t$。对该振荡积分重复分部积分 $N$ 次，每次产生一个 $\sqrt{\Delta t}$ 因子：
    \[
    \left| \int_{|\eta| \ge \varepsilon/\sqrt{\Delta t}} e^{i m |\eta|^2/2} \cdots \right| \le C_N \Delta t^{N/2}.
    \]
    因此，$|\xi| \ge \varepsilon$ 区域的贡献是 $O(\Delta t^\infty)$（即快于任何 $\Delta t$ 的幂次）。
  
  \item 有效支集限制在 $|\xi| \lesssim \sqrt{\Delta t}$ 范围内。
\end{enumerate}

\begin{lemma}
  当 $N_{\mathrm{cut}} = 2$ 时，对任意光滑函数 $f$，
  \[
  \int d\xi \, K_{\Delta t}(\xi; y) f(\xi) = \int_{|\xi| < \delta} d\xi \, K_{\Delta t}(\xi; y) f(\xi) + O(\Delta t^\infty),
  \]
  其中 $\delta > 0$ 是任意小于单射半径 $r_{\mathrm{inj}}(y)$ 的常数。
\end{lemma}

\section{对路径积分的推论}

结合以上分析，我们有：

\begin{enumerate}
  \item \textbf{一般 $N_{\mathrm{cut}}$（分布论）}：$K_{\Delta t}(\xi)$ 作为分布的支集为 $\{0\}$。它与光滑函数的配对仅依赖 $\xi=0$ 处信息。因此 $\xi = \exp_y^{-1}(x)$ 只需在 $\xi=0$ 处有定义——这总是成立的。
  
  \item \textbf{$N_{\mathrm{cut}}=2$（显式估计）}：有效支集在 $|\xi| \lesssim \sqrt{\Delta t}$ 内。对紧致 $M$，全局单射半径 $r_{\mathrm{inj}}(M) = \inf_{y\in M} r_{\mathrm{inj}}(y) > 0$。取 $\Delta t$ 使得 $\sqrt{\Delta t} < r_{\mathrm{inj}}(M)$，则 $\xi$ 的有效取值始终落在 $\exp_y$ 为微分同胚的球 $B_{r_{\mathrm{inj}}}(0)$ 内。因此 $\xi = \exp_y^{-1}(x)$ 唯一存在且光滑。
\end{enumerate}

两种情形都保证了在路径积分的每一步中，
\[
\xi_i = \exp_{x_i}^{-1}(x_{i+1})
\]
在短时传播子的有效支集上是良定义的。\end{document}
